\section{Discussion}\label{sec:honest}

\subsection{An honest account of the classical hardness}
We state plainly what is certified and what is open. Sampling from a subclass of unitary-cluster-Jastrow
circuits is classically hard---it embeds instantaneous-quantum-polynomial
computation~\cite{bremner2016,hafid2025}---and adding a single non-Clifford resource already makes
distribution learning hard~\cite{hinsche2023}, just as the difficulty of learning fermionic
states grows with their non-Gaussian gate count~\cite{melehera2025}; the cost of matchgate simulation grows with the
non-Gaussian (controlled-phase) content of the circuit~\cite{extraferm}. These are
not, however, a monotone certificate of hardness for a specific trained instance: the fermionic
AntiFlatness is a \emph{state} monotone while the matchgate simulation cost is a \emph{circuit}
quantity, and their relation is empirical rather than a theorem. A tensor-network comparison
sharpens the scope. On one-dimensional chains a matrix-product state of small bond dimension
represents these states far more compactly than any determinant subspace---unsurprising, since one
dimension is where the density-matrix renormalization group is near-optimal~\cite{schollwock2011}. DMRG has
matured into a production \emph{ab initio} method for quantum chemistry~\cite{chanheadgordon2002,block2},
and mature correction-vector, dynamical-DMRG, time-dependent-DMRG~\cite{paeckel2019}, and kernel-polynomial
spectral-function methods already
exist~\cite{kuehner1999,jeckelmann2002,holzner2011,nocera2016,whitefeiguin2004,kpm2006};
beyond tensor networks, classical Pauli-propagation methods now evaluate real-time correlators
directly~\cite{kemper2026}.
\emph{We therefore do not claim advantage over tensor networks in one dimension.} The
configuration-selection advantage we demonstrate---over deterministic Krylov and CIPSI selectors, not
over heat-bath CI---is relevant in the settings where matrix-product representations strain: genuinely
two-dimensional geometries and quantum chemistry, whose orbital graphs lack the one-dimensional locality
that matrix-product states exploit. Figure~\ref{fig:ladder} isolates the resource mechanism in a
controlled contrast: at \emph{matched} one-body magic a two-leg Hubbard ladder already carries a larger
minimal bond dimension---and a $|\mathcal S|$ up to $\sim\!2.2\times$ that of the chain, needing a
correspondingly larger $(N{\pm}1)$ fraction---the chain-to-ladder ratio growing with size. A ladder of this
width ($\chi\!\approx\!38$) remains classically near-exact for dynamical DMRG, so this is a resource
contrast---the cost tracks $\chi$, not $\mathcal F_1$, exactly as the resource theory predicts---not a
beyond-classical claim. On the positive side,
Theorem~\ref{thm:b2} (App.~\ref{app:b2}) guarantees that a subspace containing the order-$K$ Krylov
space reproduces the exact spectral moments through order $2K{+}1$ and is captured with a polynomial
shot budget, so the sampled reconstruction is convergent and its error is bounded by the captured
spectral weight.

\paragraph{The strongest classical counterexamples, and why they sharpen rather than sink the thesis.}
Two recent results set the bar for any ground-state advantage claim, and we concede both without
reservation. Reinholdt \emph{et al.}~\cite{reinholdt2025} show that quantum-selected configuration
interaction re-samples already-seen determinants in a coupon-collector process and returns expansions
\emph{less} compact than classical heat-bath or adaptive selected CI; Zhang and
Otten~\cite{zhangotten2026} reach the ground state of the $[\mathrm{4Fe\text{-}4S}]$ cluster from
\emph{random} Slater determinants using of order $10^{6}$ times fewer determinants and CPU-hours than the
quantum demonstration. For the ground-state \emph{energy} of the systems shown to date, strong classical
selectors match or beat sample-based diagonalization, usually far more cheaply---we do not dispute it.
Our decoupling explains \emph{why} this must be so rather than contradicting it. The quantity these
critiques probe is the size and compactness of the determinant expansion $|\mathcal S|$, a
basis-dependent count that a free orbital rotation moves across its whole range and that the invariant
one-body magic $\mathcal F_1$ is provably blind to (Sec.~\ref{sec:resource}, Thm.~\ref{thm:witness}). The
critiques therefore land where our analysis predicts a classical method should win---on a
basis-dependent, higher-order cost that no efficiently measurable state monotone we know can rank---and
they say nothing about the invariant multireference diagnostic, which was never a cost claim. The
conclusion is not that the critiques are wrong but that the advantage question has been posed about the
wrong observable: the ground-state energy, where $|\mathcal S|$ is classically tamable and random- or
heat-bath-selected determinants suffice, rather than the frequency-resolved response, whose faithful
reconstruction demands a large fraction of the $(N\pm1)$ sector (Sec.~\ref{sec:resource}) and where the
compactness edge that powers heat-bath and adaptive selected CI~\cite{holmeshci2016,sharmashci2017,tubmanasci2016} does not transfer. We make no beyond-classical claim
even there; we relocate the open question to where it remains open.

\subsection{Machine learning that preserves, rather than removes, the quantum sampling role}
Machine learning can enhance the pipeline only where it \emph{consumes device samples or schedules
device work}---the regime where learning from quantum experiments carries a provable
advantage~\cite{huang2022}---never where it could generate the dominant determinants unaided, which
would dequantize the method exactly as the classical selectors of the introduction
do~\cite{belagali2026,hinqs2026,review}. We put the strongest candidate to a stringent test. A
generative tail-completer---an orbital-native generative flow network~\cite{bengio2021}
trained to sample high-weight $(N{+}1)$-sector configurations---is asked to complete an under-sampled
device head, in both device-seeded and from-scratch conditions (Fig.~\ref{fig:gflow}).

The result that carries the effect size here is not the generator but the baseline it fails to clear, and we
state it in its own right rather than as a foil. Adding a cheap classical Epstein--Nesbet perturbative
(CIPSI) tail~\cite{cipsi1973} to the \emph{same} device histogram---no additional shots, no additional
quantum processing, and a first-order selection score that requires a single application of $H$ to the
sampled wavefunction---cuts the energy error from $39.24\pm3.95$ to $26.98\pm0.66$~mHa (mean$\pm$s.d.\ over five
seeds), a $31\%$ reduction and $3.1$~s.d.\ of the pre-correction spread. That $12.26$~mHa is three times
the $4.04$~mHa the generative model then contributes on top of it, and it costs no quantum resources at
all; it should be the default post-processing step on any under-sampled sample-based head; a demonstration that omits it is
comparing against a baseline that is too weak. Only against that corrected baseline is the generative
question well posed---and there the learned generator's point estimate edges lower still but not by a
statistically resolved margin (the further $4.0$~mHa fused-vs-control difference is $1.2$~s.d.\ of the
generator's own five-seed spread, against $3.1$ for the classical step; Fig.~\ref{fig:gflow}), and the residual
error is closed by drawing more device samples, not by the generator. This is a
\emph{null result}: on this single instance a learned generator does not beat the cheap classical baseline
that the quantum sampling already enables---it does not, on the evidence here, manufacture dominant
determinants the sampling misses. The safe roles for machine learning are therefore strictly
post-processing and scheduling---mitigation and certification of the sampled histogram feeding
recovery~\cite{qesem2025,algorithmiq2026}, a role in which generative flow networks are already being
deployed for measurement optimization---both to group commuting operators into compatible measurement
sets~\cite{vargashernandez2024} and to optimize the discrete measurement
basis~\cite{vargashernandez2025}---and an active-learning schedule that
places the evolution times where the spectral function is most uncertain, in the spirit of Bayesian
optimization for inverse problems in quantum reaction dynamics~\cite{vargashernandez2019}---none of which replaces the quantum
sampling. The dequantization guard reported here corroborates, from the spectral-function side, the
classical-simulability frontier mapped for ground-state SQD in our companion review~\cite{review}.

This delimits, rather than forecloses, the ``quantum-for-AI'' horizon, and bounds where it could lie. The
2026 generative frontier has two mirror faces: Born machines that are classically \emph{trainable} yet
quantum-hard to \emph{sample}~\cite{genqadv2025,spectralborn2026}, and sample-based diagonalization,
which is quantum-\emph{sampled} yet classically \emph{diagonalized}. The defensible role for a learned
model is the one that respects this duality---consuming the device histogram rather than replacing
it~\cite{huang2022,bennewitz2022}---and the theme of the present programme, amortization, names the
opening the current instance-by-instance generative selectors~\cite{nqssc2026,zeni2026,pigensqd} leave
open. Per-molecule generative configuration recovery can already improve a \emph{single} target's sampled
sector when retrained for that molecule~\cite{pigensqd}; what our negative result and this programme
target is different---a recovery model trained \emph{once} across a family to complete the
physical-sector tail of a new, under-sampled instance without per-instance retraining. We take a first,
falsifiable step on this question (Fig.~\ref{fig:amort}). Trained once across a six-instance Hubbard family
and evaluated leave-one-out, an amortized configuration-importance model recovers a held-out instance's
spectral support \emph{at the level of} the strong per-instance classical selector (an Epstein--Nesbet
tail it never sees): the leave-one-out relative-$L_1$ difference is $-0.01$ to $-0.09$ across subspace
fractions, within $\sim\!1$~s.d.\ over the six folds. Amortization without per-instance retraining is
therefore \emph{feasible}---a train-once model generalizes across the family---but shows \emph{no
statistically resolved advantage} over the classical baseline yet, the status of the open
question: the single-instance
negative result sets the bar, the amortized result meets but does not clear it, and---consistent with
both---any advantage, if it exists, would be confined to the high-noise, high-correlation regime, never on the easy instances
where the classical prior already wins.

\begin{figure}[htbp]\centering\includegraphics[width=0.80\textwidth]{figs/fig_gflownet.pdf}
\caption{\textbf{The large effect is the free classical tail; the learned generator adds no resolved
gain.} At an under-sampled device head ($\mathrm{N_2}$, $1000$ shots, $5$ seeds), a cheap classical
Epstein--Nesbet (CIPSI) tail~\cite{cipsi1973}, drawing no additional shots, cuts the error from
$39.24\pm3.95$ to $26.98\pm0.66$~mHa---a $31\%$ reduction and $3.1$~s.d.\ of the pre-correction spread.
A generative GFlowNet tail-completer then edges the point estimate lower still ($25.96\pm3.40$~mHa from
scratch, $22.94\pm3.30$~mHa fused onto the classical prior), but that further $4.0$~mHa fused-vs-control
improvement---a third the size of the classical step---is only $1.2$~s.d.\ over five seeds, not a
statistically resolved win. The residual error is closed by drawing more device samples, not by the generator: on this
single instance the quantum sampling, not a learned generator, supplies the dominant determinants that
additional device shots recover.}\label{fig:gflow}\end{figure}

\begin{figure}[htbp]\centering% G2 amortized recovery (fig:amort) -- amortized_recovery.json. Train-once-across-family, leave-one-out over 6 Hubbard instances.
\definecolor{amGen}{HTML}{D55E00}\definecolor{amCls}{HTML}{0072B2}\definecolor{amInk}{HTML}{1B1B1B}
\begin{tikzpicture}
\begin{axis}[width=8.6cm,height=6.2cm,paperaxis,
  xlabel={subspace fraction $|\mathcal S|/D$}, ylabel={spectral relative-$L_1$ (leave-one-out)},
  xmin=0.02,xmax=0.33, ymin=0.35,ymax=1.05, xtick={0.05,0.1,0.2,0.3}, ymajorgrids, grid style={draw=black!10},
  legend style={font=\footnotesize,at={(0.5,1.02)},anchor=south,legend columns=2,/tikz/every even column/.append style={column sep=12pt},draw=none,fill=none},legend cell align=left,
  error bars/y dir=both,error bars/y explicit]
\addplot[amCls,line width=1.5pt,mark=square*,mark size=2.8pt,mark options={fill=amCls,draw=white}] coordinates {(0.05,0.877) +- (0,0.101) (0.1,0.856) +- (0,0.114) (0.2,0.776) +- (0,0.100) (0.3,0.657) +- (0,0.093)}; \addlegendentry{classical (per-instance)}
\addplot[amGen,line width=1.5pt,mark=*,mark size=2.8pt,mark options={fill=amGen,draw=white}] coordinates {(0.05,0.869) +- (0,0.161) (0.1,0.833) +- (0,0.162) (0.2,0.688) +- (0,0.216) (0.3,0.618) +- (0,0.199)}; \addlegendentry{amortized (train-once)}
% (annotation removed to avoid overlapping the data/error bars; the caption states the bands overlap)
\end{axis}
\end{tikzpicture}

\caption{\textbf{Amortized configuration recovery: a first data point on the programme's open question.}
A configuration-importance model trained \emph{once} across a six-instance Hubbard family
($L=6,8$; $U/t=4,8,12$) and evaluated \emph{leave-one-out} recovers a held-out instance's
$(N{+}1)$-sector spectral support, versus the strong per-instance classical selector (Epstein--Nesbet,
untrained), as a function of the sampled subspace fraction. The amortized model---which never sees the
held-out instance---meets the classical baseline (leave-one-out $\Delta=$ amortized $-$ classical
$=-0.01$ to $-0.09$ across fractions, within $\sim\!1$~s.d.\ over the six folds): amortization without
per-instance retraining is \emph{feasible} but shows no statistically resolved advantage yet. Error bars
are $\pm1$~s.d.\ across the six leave-one-out folds.}\label{fig:amort}\end{figure}

\subsection{Significance for academia and industry}

\paragraph{Academia.} The result reframes the quantum-advantage question in electronic structure. The
debate has centred on the ground-state \emph{energy}, where the evidence for a generic exponential
advantage is weak~\cite{leechan2023} and the flagship sample-based demonstrations are being
dequantized~\cite{belagali2026,reinholdt2025}. We move the question to a regime where an advantage is
defensible---dynamical response---and, more consequentially, replace an asserted advantage with a
\emph{precise resource statement}. One-body fermionic magic, though efficiently computable, is
decoupled from the sampling cost and cannot certify where a quantum processor is required; what it
faithfully reports is the static, multireference character of a target, while the operative hardness
lives in the higher-order, connected-cumulant sector. In drawing this line we tie the resource theory
of fermionic non-Gaussianity~\cite{hebenstreit2019,sierant2026,tarabunga2026} to sample-based
electronic-structure computation~\cite{robledomoreno2025,kanno2023} for the first time, and open a
concrete programme: the frequency-resolved observables now accessible from sampled subspaces are the
very objects of dynamical mean-field theory~\cite{dmft1996} and correlated
spectroscopy~\cite{damascelli2003}, and the resource analysis is a template for how quantum-advantage
claims can be \emph{scoped and delimited} rather than asserted---a discipline the field increasingly
demands~\cite{leechan2023,babbush2021,cerezo2025bp}.

\paragraph{Industry.} The significance for the quantum-computing industry is threefold and concrete.
\emph{First, the right observables.} What the method produces---the single-particle spectral function
measured by photoemission, its momentum-resolved form, and the dynamical structure factor probed by
neutron and X-ray scattering~\cite{squires2012,ament2011}---are precisely the quantities that
materials discovery, catalysis, and correlated-electron device design require, and that the
energy-only output of today's sample-based hardware does not
deliver~\cite{robledomoreno2025,extsqd2025}. \emph{Second, a correlation diagnostic.} Quantum runtime on the quantum-centric
supercomputers now in operation---an IBM Heron processor coupled to the Fugaku
supercomputer~\cite{robledomoreno2025,shirakawa2025}, the QSCI software
stack~\cite{kanno2023,adaptqsci}, and the trapped-ion Green's-function and
dynamical-structure-factor pipeline~\cite{greenediniz2024,quantinuum_dsf2026}---is scarce and
expensive, so knowing \emph{which} targets are strongly correlated is operationally valuable. The
efficiently measurable fermionic magic flags exactly the multireference regime where mean-field and
weakly-correlated classical methods break down. It is not, however, a cost oracle: being
decoupled from the determinant support (Sec.~\ref{sec:resource}), it cannot by itself rank a strongly-correlated target against a
classical selected configuration interaction~\cite{sharmashci2017,tubmanasci2016}, the density-matrix
renormalization group~\cite{whitedmrg1992}, or auxiliary-field quantum Monte
Carlo~\cite{mottazhang2018}---that decision turns on the higher-order, basis-dependent cost the
one-body magic does not see~\cite{reinholdt2025}. \emph{Third, low adoption cost.} The method is
hardware-native: it needs only computational-basis sampling, with no Hadamard tests or controlled
unitaries~\cite{bauer2016,endo2020}, so it runs on current devices within the existing SQD stack and
its GPU-accelerated classical back end~\cite{gpusqdomp,gpusqdthrust}, and it \emph{lightens} rather than adds to the
observable-estimation and error-mitigation burden that dominates near-term
cost~\cite{qesem2025,algorithmiq2026,qedma2026}. Fermionic magic is thus offered as an efficiently
measurable multireference diagnostic---a \emph{one-sided} triage that flags, before committing runtime,
the strongly-correlated targets on which mean-field and weakly-correlated classical methods fail. It is
expressly \emph{not} a placement oracle: whether a workload belongs on the quantum processor turns on the
basis-dependent cost---the bond dimension $\chi$---that one-body magic provably cannot supply
(Sec.~\ref{sec:resource}), and we claim no scheduling result absent a scheduling experiment.
